% Nouveau EGA V (cote 25): archive cote 25, site page 149 (PDF page 150).
% English translation of the checked French annotated master.
% AI-translated by GPT-6 Astra (OpenAI, ChatGPT Pro mode).
\scanpage{25/149}
\mspage{\hand{10}}
\var{\margin{\unsure{$\diagdown$}}}{\placedmargin{\unsure{$\diagdown$}}}{marginal note placed}{ega25_25_149-v0001}
% The long diagonal source stroke has uncertain function; the text remains live.

Having reached this admirable degree of generality, I unfortunately
realize that these developments are still too particular, and that instead
of a subprescheme defined by a single section of $\underline{O}_X$, one should
consider the one defined by m sections, or, what\hand{/}amounts to the same thing, by a
homomorphism $\underline{E} \hand{\longrightarrow} \underline{O}_X$, where $\underline{E}$ is a loc free Module of rank m. I
leave it to you to decide whether to begin from the outset with the
general case, following the principle that one must proceed from the general to the
particular.

Let us first introduce terminology for rings A with
augmentation $A \hand{\to} A/I$ (which will also be used in the sheaf
context). The augmentation is called \uline{regular} (or should we say
quasi-regular, with the definitions adopted in section 19?) if \var{\inserted[typed]{a)}}{a)}{insertion}{ega25_25_149-v0002} $I/I^2$ is
a projective A/I-Module of finite type, and \var{\inserted[typed]{b)}}{b)}{insertion}{ega25_25_149-v0003} the homomorphism
\[ \theta\colon \mathrm{Sym}\,(I/I^2) \hand{\longrightarrow} \mathrm{gr}_I(A) \]
is an isomorphism. We shall say that the augmentation is \var{\inserted{of type}}{of type}{insertion}{ega25_25_149-v0004} ``\uline{ordinary}
\uline{quadratic}'' when a) above holds, as well as b') the preceding canonical
homomorphism has a kernel generated by $\theta^2$, which kernel $\underline{K}$
is a ``nondegenerate'' invertible \var{\inserted[typed]{sub-}}{sub-}{insertion}{ega25_25_149-v0005}Module of $\mathrm{Sym}^2(I/I^2)$. By ``nondegenerate''
we mean that the corresponding homomorphism
\[ \underline{K} \otimes (\underline{I}/\underline{I}^2)^{\hand{\vee}} \hand{\longrightarrow} \underline{I}/\underline{I}^2 \]\var{\margin{\unsure{$\prime$}}}{\placedmargin{\unsure{$\prime$}}}{marginal note placed}{ega25_25_149-v0006}
(obtained using the canonical homomorphism
\[ \mathrm{Sym}^2(I/I^2) \hand{\longrightarrow} \mathrm{Hom}\ (\ (I/I^2)^{\hand{\vee}}\ ,\ I/I^2\ )\quad ) \]
is bijective, or equivalently surjective (since these are
projective Modules \var{\inserted[typed]{over A/I}}{over A/I}{insertion}{ega25_25_149-v0007} of the same rank at every point), or again: taking
locally a basis \var{\inserted{$\phi$}}{$\phi$}{insertion}{ega25_25_149-v0008} of $\underline{K}$, \var{\xout{the corresponding homomorphism}}{}{typewriter strike-out}{ega25_25_149-v0009} the corresponding
quadratic form on $(I/I^2)^{\hand{\vee}}$ is nondegenerate. It suffices

\mspage{\hand{149}}
\tnote{25-149-1}{The long diagonal stroke remains of uncertain function and does not remove the running text. Loc abbreviates locally. The source admitable is translated as admirable without imitating its spelling. The phrase ``kernel generated by'' followed by the squared theta is retained as written; no different kernel expression is substituted.}
